\documentclass[../../report.tex]{subfiles}
\begin{document}
\section{Data Preparation}

In order to train the models, a melody dataset is required. As discussed before,
this can be extracted from MIDI files. This project uses the Lakh MIDI dataset
\cite{Raffel2016}, more precisely its `clean MIDI subset'. It contains
\num{17244} musical compositions of various genres.

\subsection{MIDI processing pipeline}

The MIDI format is too verbose to be fed directly into the model. Using
Magenta's \texttt{pipelines} package, the MIDI files are taken through several
preprocessing stages to produce appropriately structured input data.

\subsubsection{Conversion}

First of all, each MIDI file is converted into a \texttt{NoteSequence} protocol
buffer\footnotemark{}, which is provided by the \texttt{note-seq} library
\cite{NoteSeq}. This library simplifies manipulation of note sequences, and its
utilities are used to implement most of the stages in this pipeline.

\footnotetext{A protocol buffer is a compact and efficient format for
serialising structured data, roughly equivalent to XML or JSON.}

\subsubsection{Partitioning}

As with all machine learning experiments, the dataset is split into training and
validation partitions. The sequences in the training set will be used to adjust
model parameters, whereas the validation set will serve as a test of the model's
predictive capabilities.

\subsubsection{Handling time changes}

MIDI files can encode music with changes in tempo or time signature. To reduce
the complexity of data that the model has to generalise, sequences with variable
timing are split into subsequences of constant timing.

\subsubsection{Quantisation}

Basic RNN assumes that melody events are evenly spaced in time. In reality, MIDI
events are bounded by arbitrary start and end times. Therefore, this stage of
the pipeline aligns the MIDI events to a fixed grid.

\subsubsection{Melody extraction}

The preceding stages apply transformations at the scale of whole MIDI files,
which may contain multiple tracks. In this stage, each track is processed
individually and monophonic melodies are extracted. There may be several
melodies produced from a single track, as segments of silence are removed.

\subsubsection{Validation}

Each melody is validated against two criteria:

\begin{itemize}
  \item Must be at least 7 bars long. Short melodies are not very useful, as
  they do not help the model learn long-term structure.

  \item Must contain at least 5 unique pitches. This ensures that the model is
  not trained on overly repetitive melodies, such as those in accompaniment
  tracks.
\end{itemize}

Melodies failing to satisfy these criteria are discarded.

\subsubsection{Truncation}

Longer melodies increase the amount of RAM required to train the model. To avoid
exceeding the hardware limitations of this project, all sequences are truncated
to 128 steps.

\subsubsection{Model encoding}

Finally, melodies are converted into the model-specific format. For Basic RNN,
this simply means mapping melody event labels to their corresponding one-hot
vectors. Let's consider the melody given in figure \ref{fig:encoding-melody}. It
could be represented as a sequence \(\bm{m}\) and then encoded as follows:

\begin{figure}[h]
  \centering
  \includegraphics[width=0.75\textwidth]{encoding-melody}
  \caption{Melody used in the encoding example.}
  \label{fig:encoding-melody}
\end{figure}

\begin{align*}
  \bm{m} &=
  \begin{bmatrix}
    0 & -2 & -2 & 1 & -1 & 2 & 2 & -2
  \end{bmatrix}
  \\
  \mathrm{encode}(\bm{m}) &=
  \begin{bmatrix}
    0 & 1 & 1 & 0 & 0 & 0 & 0 & 1 \\
    0 & 0 & 0 & 0 & 1 & 0 & 0 & 0 \\
    1 & 0 & 0 & 0 & 0 & 0 & 0 & 0 \\
    0 & 0 & 0 & 1 & 0 & 0 & 0 & 0 \\
    0 & 0 & 0 & 0 & 0 & 1 & 1 & 0 \\
  \end{bmatrix}
\end{align*}

The above encoding assumes eighth note time steps. Three possible \emph{note-on}
events are encoded by values 0, 1 and 2. Special value \num{-2} indicates
\emph{no-event}, i.e. continue playing previous note; \num{-1} represents a
\emph{note-off} event.

The real model uses sixteenth note time steps. Also, the one-hot vectors have 38
entries, representing 36 \emph{note-on} events and the two special events.

\end{document}
